\section{Tier 1: Validating the inverse-MC estimator inside the filter}
\label{sec:phase1-synthetic}

This tier is a code and mathematics gate, not yet a sat-navigation
detection experiment. It asks one narrower question: if the true Markov
chain is known synthetically and only a subset of its states is observed,
does this implementation recover the stationary distribution at the
unobserved states about as well as Morimura et al.'s published method?

The reproduction follows Morimura et al.'s synthetic benchmark on a
random $n=100$ softmax-parametric chain and matches the published
proposed-method RMAE curve to within $\sim 0.05$ for $|\Xo|\ge 10$.
RMAE is the relative mean absolute error of the recovered stationary mass
on unobserved states; lower is better. The only material deviation is the
$|\Xo|=5$ setting, where the $75/25$ validation split leaves roughly one
validation state and makes $\lambda$ selection degenerate.

The synthetic reproduction is therefore sufficient as an implementation
gate for the downstream tiers: it verifies the analytic gradients, the
partial-observation loss, and the estimator's stationary-distribution
recovery under the benchmark's own assumptions. It says nothing by itself
about the non-stationary SUMO and Xuancheng windows studied later. Full
protocol details, the RMAE table, and the mathematical justification for
the scale-free prediction convention are given in
Section~\ref{app:phase1-synthetic-details} in the Appendix.
